\section{The Two Cosmic Stories}
\label{app:cosmic-stories}

The oldest question in cosmology is whether the universe \emph{began} or has \emph{always been}. Vibe theory does
not answer it. The substrate admits two accounts that are equally self-consistent and equally complete, and the
fair thing to do is to set them side by side and let the reader weigh them. We tell both in equal measure. This
chapter is the conceptual centrepiece of the part, a fork offered as a gift rather than a verdict. The diptych
comes first, then each story in full, then the perennial images that older traditions reached for again and again.

\subsection{The shared ground}

Before the fork there is common ground, and it is worth naming because it is the part most people find hardest to
accept. There was never a \emph{nothing}. Neither story is a story of something arising out of absolute void. The
most uncomfortable idea of all, a bare nothing that somehow produces a world, appears in neither picture. Every
ancient image agrees on this. There is always a prior fullness, a fire, an ocean, the One, a dreamer, a breath.
The world is a spark of the fire, a wave of the ocean, a note of the song.

The reason to believe this is stronger than taste. Take the word nothing strictly, not empty space, which already
has size and directions and a clock, but the total absence of any distinction at all. A state like that cannot
hold, for four plain reasons. First, it is a balance point and not a floor. A uniform state with no distinctions
is like a ball on the very top of a round hill, and the slightest difference rolls it off because there is no
restoring force, where a floor would catch it. Second, no difference is itself a difference. The moment the state
can be named as one whole, distinct from what it is not, a line has been drawn, so a perfectly seamless unity
cannot even describe itself without breaking its own seam. Third, a whole that is everything has nothing else to
relate to, and since to exist is at bottom to relate, it can only relate to itself, and self-relation splits it
into a noting end and a noted end. That split is the first distinction. Fourth, there is no road back, because
erasing a distinction is itself an event, a new distinction between before and after, so every attempt to reach
nothing produces more something. Nothing is unstable, self-contradictory, forced to turn on itself, and one-way.
It is the one state reality cannot be in.

This is why something is the default and nothing is the thing that needs explaining, the reverse of the usual
intuition. It is also where the two stories take their cleanest forms. Read the self-relation as a first event and
you have Story A, the One waking, a forced first split. Read it as what has always already been happening, with no
first, and you have Story B. The arguments against nothing are shared. Whether they imply a true first distinction
or an eternal one is exactly the fork, and we leave it open.

The base is discrete and computable, and computability is the referee that lets both stories stand. As the next
chapter sets out, computability forbids the continuum and supertasks. It does not forbid an infinite past. So a
beginningless history is no less computable than a finite one. Both stay within what an idealized machine could in
principle run.

Our observed big bang fits both readings without strain. It can be the first morning of everything, the true
edge of time. It can also be a single spark struck in a vast and eternal flame. Nothing in the local data of our
chapter of cosmic history settles which. Our visible universe is not necessarily the whole, and the question of
the whole is exactly the question we are leaving open.

\subsection{The mirror symmetry}

The two stories are mirror images of one another, and seeing the symmetry is most of the point. One has a first
point that is a clean edge, like the North Pole of a sphere. The other has no first point at all, like a line that
runs in both directions, or the negative integers with no smallest member. Both refuse creation from nothing. The
whole of cosmology's deepest question reduces to a single fork. Is there a north-pole edge to time, or is there
not.

\begin{table}[ht]
\centering
\begin{tabular}{@{}lll@{}}
\toprule
 & Story A, a beginning & Story B, beginningless \\
\midrule
the ``before'' question & a real edge, no before & no edge, always a before \\
what the universe is & the One unfolding into the many & a verb, eternally becoming \\
the engine & self-noticing instability & self-sustaining churn \\
something from nothing & no, there was the One & no, there was always becoming \\
our big bang & the first morning & a spark in the eternal flame \\
\bottomrule
\end{tabular}
\caption{The two origin stories as mirror images. Each refuses creation from nothing.}
\end{table}

\subsection{Story A, the One unfolding}

This is the account in which the universe truly began. In it the question \emph{what started the universe} is not
malformed. It has an answer. The universe started when the One first noticed itself. Before that there was not
nothing. There was unity, the whole of what-is, seamless and undivided. The beginning is the moment that seamless
one first split.

Picture a single awareness with no parts, no inside or outside, just one. Here is the engine. To be is to be
aware, and to be aware is to be aware of something. The One has nothing else to be aware of, so the only thing it
can notice is itself. The instant it notices itself there are two ends, the noticer and the noticed. Awareness of
self is already a split. The universe is the One waking up, and waking up is becoming two. The first beat is that
first noticing. It is not a thing made from nothing. It is the One differentiating into the many.

The fracture was forced, not accidental. A structure rich enough to notice itself cannot hold a single consistent
state, because the self-as-noticed always differs from the self-as-noticer. This is a clean logical fact, of the
same family as the diagonal argument. Unity is then like a ball balanced exactly on the tip of a peak. The balance
point is real, but the slightest disturbance sends it rolling, and there is always a slightest disturbance. Unity
had to fall into difference. The beginning is not a matter of why it started. It is that it could not avoid
starting.

The first moment has no before, and that is not a hole. We feel that a first moment is impossible, because we keep
asking what came before it. The clean answer is that \emph{before} is a relation inside time, and time itself
begins at the first moment. So there is no before, for the same reason there is nothing north of the North Pole.
The North Pole is a real, definite, northernmost point. Ask what is north of it and there is no answer, not
because something is hidden but because northness ends there. The first moment is the North Pole of time, and the
missing before is not a gap. It is the edge itself.

Once there are two, there are relations between them to notice, and noticing those makes more, and noticing those
makes more again. It runs away, like a microphone held to its own speaker, each output becoming new input until
the room is full of sound. The first tiny split feeds the roaring, branching, growing structure that is the
universe. From one forced fracture the whole mesh unfolds, beat by beat, the many blooming out of the one.

The finite seed is also the natural place to start thinking, even for a reader who ends up an eternalist. Assuming
a first seed makes the picture concrete. You can watch the universe build itself from nothing but the rule, step
by step, and only then ask whether the first step was truly first. Run forward from the seed and a clean mechanism
appears, which we may call birth at peace. Each beat, reflection adds a new ring of cells at the growing frontier,
and every newborn cell starts at tone $0$, peace. This is forced twice over. Conservation forbids a nonzero birth,
since a $+1$ or $-1$ cell would add charge and break the balance, and peace is the natural blank value anyway. The
newborn then acquires its tone from the world, as the conserved charge of the existing crystal flows into it along
its connections, zero-sum, so the total never changes. The frontier is a blank peaceful canvas, and the world
paints it with charge it already had. The genuine freedom at the frontier is freedom of relation and of
redistribution, never freedom to mint charge. This same mechanism works identically on the eternal reading, where
every newborn shell is born at peace whether or not there was a first one.

No single image carries the story alone, because each one smuggles in a prior context. The One waking up suggests
a prior sleeper, which the North Pole image removes by denying any before. The ball on the peak seems to assume a
ball and a gravity set up in advance, which dissolves once it is read as the logical instability of
self-reference rather than a mechanical picture. The North Pole image itself has no hole, being a clean geometric
fact about a sphere. The One unfolding into the many makes plain that this is not something from nothing, since
there was always the One. Fused into one sentence, the story runs like this. Unity is a ball on the tip of a peak
that cannot help but roll, the One waking into two, and the moment it wakes is the North Pole of time, with no
before because before-ness begins there, and from that one forced fall the many bloom forever.

\subsection{Story B, the eternal becoming}

This is the account in which the universe never began. The question \emph{what started the universe} is still an
excellent question, and following it carefully leads somewhere subtle. Starting is something that happens inside
time. A song starts, a race starts, a fire starts, and every start you have ever seen had a before, a moment of
quiet before the song, a held breath before the race. Start means something only when there is an earlier time for
the starting to happen in. To ask what started time itself is to ask for a before-the-first-before, and that is
where the question quietly looks for an earlier time when there may be none. The deepest answer is not a clever
first cause. It is that time may not need a starting line at all, because a starting line is itself an in-time
idea.

The universe is a verb, not a noun. We are used to things being made and then sitting there existing, a table
built, a star born. So we imagine the universe the same way, made once and then existing. But the universe is not
a table. It is a happening, like a flame, like a river, like a song being sung right now. A flame is not made and
then existing. It is continuously burning. Stop the burning and there is no flame. A verb does not need a maker.
It needs only to keep happening.

Here is the heart of it. In this universe nothing is ever lost, because the rule runs backward as perfectly as
forward. Every now carries its before inside it, fully. The present is the past transformed, not erased. Take any
moment and ask what came before, and there is always a real, definite answer, because the moment literally
remembers it. That moment remembers its before, and so on, with no bottom, the way the counting numbers have no
smallest one. There is never a now that came from nothing, because every now is the transformation of a real
before.

On this reading the strange option is a beginning, not the lack of one. Look at what a real beginning would
require. At a single instant the entire universe leaps out of absolute nothing, uncaused, for no prior reason, and
then keeps going. That is the gut-wrenching, impossible-feeling thing, something from nothing, a causeless flash.
The beginningless picture is gentler. Nothing ever came from nothing, because there was never a nothing to come
from. Existence did not have to be switched on.

What, then, is the big bang. Our visible universe really does seem to begin about $14$ billion years ago. But that
need not be the beginning of everything. It can be the beginning of our chapter, a bubble forming, a new region
catching fire in the endless burning. Page one of our story, not page one of the book. There were always pages
before ours, and the book has no first page. Both are true at once. Our local universe began, and the whole was
always going. The big bang is a sunrise, not the creation of the sun.

Three metaphors plug each other's holes here. The eternal flame says that the universe is a happening, not a
thing, but real flames are lit and go out. Counting backward says there is no first and no one finds that strange,
which plugs the flame's was-lit hole, but numbers are static. The reversible film says every moment is the exact
reshaping of the one before, nothing lost, a film that plays the same backward as forward so that from any frame
both neighbours can be rebuilt. This is literally true of the rule, and it is also the mirror image moved from
frozen space into flowing time. Fused into one sentence, the universe is an eternal flame whose every flicker is
the exact reshaping of the flicker before it, nothing ever lost, with no first flicker, the way counting backward
has no smallest number.

The reversible film is not only a metaphor. It was tested. Start from random noise rather than a hand-picked
state, run the rule forward $50$ steps, then backward $50$, and the exact starting state returns with zero error.
The activity holds steady, around two-thirds of cells changing each tick, for hundreds of ticks. It never freezes
and never blows up. The decisive part is the control. An irreversible blurring rule from the same random start
decays to a frozen dead state within about $200$ ticks. A reversible churn can run forever with no start and no
end. An irreversible one cannot. This shows the eternal picture is fully self-consistent. It does not prove the
universe is eternal, only that it is allowed \cite{pollard2026vibetest}.

One real subtlety remains. Because the rule is time-symmetric, the arrow of time cannot come from the rule
itself. It must come from the boundary or the initial conditions, the so-called past hypothesis. The reversible
vibe rule inherits this famous feature, so we do not claim to have produced time's arrow out of nothing. There is
a clarifying payoff in the growth picture all the same. The single growth of one new shell per beat is at once the
time step, the arrow, and the spatial expansion, three usually separate things made one. Cosmic expansion is not
inserted by hand. Time is the growing of space, in the form of a de Sitter law $a(t) \sim e^{Ht}$ with $H \approx
0.80$ per beat. This holds on either origin reading.

The beginningless story stays computable because at any flicker only finitely much has burned. There is always
more before and after, but never a completed whole. The history is endless but stepwise, never grasped all at
once. So an infinite past is allowed while the continuum and supertasks are not.

\subsection{The perennial images}

Humans have always reached for the same images when they picture the local world as a part of a vast or eternal
whole, and the One becoming the many. These pictures resonate because they are old and deep, and the structures
they describe turn out to match what the substrate produces. The deepest shared point in all of them is that none
is creation from nothing. Each carries a prior fullness, and the world is a part of it.

The fire images run across both stories. A spark in the flame is the ever-living fire of Heraclitus, the Stoic
logos, the divine spark of the Gnostics and of Eckhart, and it serves the beginningless reading. A flame lit from
a flame, one candle lighting another undiminished, is the Buddhist image for passing on without loss, and it
serves both, since the new fire is real and begins yet nothing is created. The first struck spark, the chord
blooming from a single note, serves the beginning.

The water images run mostly to the eternal whole. A wave on the ocean is the Vedanta image, each wave rising and
falling yet being the ocean. A bubble or foam on the sea is the Diamond Sutra's dream, bubble, dewdrop, and flash
of lightning. A drop returning to the ocean is the dewdrop of Rumi slipping into the shining sea. A whirlpool in
the river is the self as a pattern in the flow rather than a substance, a standing eddy the water passes through.
Ripples from a stone, the first disturbance spreading, serve the beginning.

The light images divide cleanly. A ray of the one sun is the emanation of Plotinus, the One radiating outward
without diminishing, and it serves both, since emanation has a source yet the source is eternal. The moon
reflected in countless waters is the Zen image of one moon and numberless whole reflections, the one appearing as
the many without splitting. A sunbeam through a window is a shaft of the one light made local.

The mind and dream images carry the same double sense. The dream of Brahma is the cosmos as the dreaming of the
divine, worlds arising and dissolving as a dreamer dreams, serving both since the dream begins and ends while the
dreamer is eternal. Indra's Net is an infinite net with a jewel at every knot, each reflecting all the others, the
part containing the whole, and this is also the literal mutually reflecting structure of the bulk. A thought in
the cosmic mind is each world arising in an endless mind.

The remaining families fill out the picture. The breath of Brahma is the universe breathing in and out across
vast ages, beginningless and cyclic, with each exhale a local beginning. The seed becoming the tree is the
Upanishadic acorn and oak, the One containing and unfolding into the many. A note in the eternal song is the
Pythagorean music of the spheres. The Word, the Om, the Logos, is the cosmos as an utterance, a first sounding
from an eternal speaker. From number and geometry come three precise images. A number in the endless count has no
smallest, the clean picture of a beginningless past. A point on an endless line has no first point. The North Pole
of the sphere is the clean picture of a first moment with no before. From weaving comes the thread in an endless
woven cloth, the Norns at the loom of time, each life a thread, which is also the literal fabric image of a self
as a braid.

The beginningless story is served by the eternal-whole images, the wave and bubble and drop in the ocean, the
spark in the flame, the number with no smallest, Indra's net, the breath of Brahma. The beginning story is served
by the unfolding images, the seed becoming the tree, the One radiating, the Word sounding, the North Pole of time.
The shared and deepest point in all of them is that neither is creation from nothing. Every ancient image has a
prior fullness, the world a spark of the fire, a wave of the ocean, a note of the song. That is what the theory
recovers. Existence is a fullness expressing itself, never a void producing a thing.

\subsection{Why we do not decide}

Nothing currently breaks the tie. Computability allows both. The reversible rule forces neither. The bootstrap
experiment shows the beginningless churn is fully self-consistent, while the finite-seed account derives a forced
first fracture, and neither is ruled out. Observation fits both. This is a genuine open question about the deepest
structure of time, and forcing a verdict would be less faithful and less interesting than presenting the fork
cleanly. The same carrying fire handles both. Was the fire always burning, or did it first flare. The fire is the
same fire either way.
